\section{Intrinsic branches on genuine normal-crossing chart overlaps}

Normal-crossing coordinate descriptions need not share a global branch
labelling. Their local branches nevertheless match uniquely. The argument
uses actual holomorphic germ factoriality and the radical Nullstellensatz,
as in \cite{demailly-cadg}, Chapter~II. It compares the original coordinate
functions themselves, without assuming a correspondence of their branches.

\begin{lemma}[A first-order holomorphic equation is prime]
\label{hir:first-order-equation-prime}
\uses{lag:germ-domain,ana:actual-holomorphic-germ-factoriality}
\lean{Hironaka.HironakaNormalCrossingsBranchOverlap.prime_germ_of_fderiv_ne_zero}
\leanok
If a holomorphic germ $f$ at $0$ vanishes at $0$ and has nonzero complex
differential there, then its genuine holomorphic germ is prime.
\end{lemma}
\begin{proof}
The germ is not a unit because its value is zero. If $f=ab$ and both
factors are nonunits, their values at zero are both zero. Choose actual
holomorphic representatives of the factors. Their product equals $f$ on
a neighbourhood, and the product rule gives $df_0=a(0)db_0+b(0)da_0=0$,
a contradiction. Thus the germ is irreducible. Actual factoriality of
the holomorphic germ ring makes every irreducible germ prime.
\end{proof}

\begin{lemma}[Actual coordinates have distinct prime branches]
\label{hir:actual-coordinate-distinct-prime-branches}
\uses{hir:first-order-equation-prime}
\lean{Hironaka.HironakaNormalCrossingsBranchOverlap.surjective_fderiv_shifted,Hironaka.HironakaNormalCrossingsBranchOverlap.prime_coordinate_of_surjective_fderiv,Hironaka.HironakaNormalCrossingsBranchOverlap.eq_of_associated_coordinates_of_surjective_fderiv}
\leanok
Let $F$ be an actual local holomorphic equivalence centred at zero, with
$F(0)=0$. Every coordinate function $F_i$ is a prime holomorphic germ,
and germs for distinct coordinates are not associated.
\end{lemma}
\begin{proof}
Differentiate the actual identity $F\circ F^{-1}=\mathrm{id}$ near zero.
It makes $dF_0$ surjective. Each coordinate row is nonzero, since a
preimage of the $i$th coordinate direction has $i$th image coordinate
one. The preceding lemma proves primeness. If $F_i=uF_j$ for a unit
germ $u$, choose its actual holomorphic representative. Differentiating
at zero gives $d(F_i)_0=u(0)d(F_j)_0$, since $F_j(0)=0$. Applying this
identity to a preimage under $dF_0$ of the $i$th coordinate direction
gives $1=0$ when $i\ne j$. This excludes association.
\end{proof}

\begin{lemma}[Equal zero unions produce a unique branch correspondence]
\label{hir:equal-zero-unions-prime-branch-correspondence}
\uses{hir:actual-coordinate-distinct-prime-branches,lag:prime-nullstellensatz}
\lean{Hironaka.HironakaNormalCrossingsBranchOverlap.exists_unique_branch_equiv}
\leanok
Let $(f_i)_{i\in I}$ and $(g_j)_{j\in J}$ be finite actual holomorphic
equations whose germs are prime and pairwise nonassociated within each
family. If their zero unions agree as germs at zero, there is a unique
bijection $\sigma:I\to J$ for which $f_i$ and $g_{\sigma(i)}$ are
associated.
\end{lemma}
\begin{proof}
Fix $i$. The product $G=\prod_{j\in J}g_j$ vanishes on the zero germ of
$f_i$, because this germ lies in the first union and hence in the second.
The radical Nullstellensatz places the germ of $G$ in the radical of
$(f_i)$. Since $f_i$ is prime, its principal ideal is prime and hence
radical, so $f_i$ divides $G$. Primeness selects a factor $g_j$
divisible by $f_i$.
Irreducibility of both germs makes them associated. Two possible $j$'s
would be associated to one another, contradicting the second family's
distinctness. Thus every $i$ has a unique match. Reversing the two families
gives a unique match for every $j$. The two match maps are inverse by the
uniqueness just proved, so the first is a bijection, uniquely determined
by association. If one index set is
empty, its zero union is empty; the same primeness argument prohibits
any branch in the other family, and the resulting bijection is empty.
\end{proof}

\begin{lemma}[Associated original equations have actual holomorphic units]
\label{hir:associated-original-equations-actual-unit}
\uses{lag:germ-domain}
\lean{Hironaka.HironakaNormalCrossingsBranchOverlap.exists_unit_equation_of_associated}
\leanok
If two actual holomorphic equations $f,g$ have associated germs at zero,
there is an actual holomorphic function $u$, nonzero at zero, such that
$f=ug$ on a neighbourhood of zero.
\end{lemma}
\begin{proof}
Association supplies a unit of the actual holomorphic germ ring with
$f=ug$ as germs. Choose an actual holomorphic representative of that
unit. Evaluation at zero is nonzero because the germ is a unit. Equality
of the actual germs is equality of their representatives on a
neighbourhood, and gives the asserted original equation there.
\end{proof}

\begin{theorem}[Actual normal-crossing flattenings agree branch by branch]
\label{hir:actual-normal-crossing-chart-branch-correspondence}
\uses{hir:equal-zero-unions-prime-branch-correspondence,hir:associated-original-equations-actual-unit}
\lean{Hironaka.HironakaNormalCrossingsBranchOverlap.exists_unique_coordinate_branch_equiv,Hironaka.HironakaNormalCrossingsBranchOverlap.exists_coordinate_unit_branch_equiv}
\leanok
Let $F$ and $G$ be two actual centred local holomorphic equivalences in
$\mathbb C^n$, both mapping zero to zero. Let $S,T\subseteq\{1,\ldots,n\}$
be finite sets of distinct coordinates. If
\[
 \bigcup_{i\in S}\{F_i=0\}=\bigcup_{j\in T}\{G_j=0\}
\]
as actual germs at zero, their original branch germs have a uniquely
produced bijection. Corresponding coordinate functions differ by actual
holomorphic units on one common neighbourhood of zero, and $|S|=|T|$.
The same assertion holds for maps centred at any common point $q$, with
equations $F_i(z+q)$ and $G_j(z+q)$ at zero.
\end{theorem}
\begin{proof}
The actual local inverse identities produce prime, pairwise distinct
coordinate germs in both families by the preceding coordinate lemma.
Apply the zero-union correspondence theorem to the selected coordinates.
Its internally produced bijection matches the entire original selected
branch germs. Apply the preceding original-unit lemma to each pair;
intersect the finitely many neighbourhoods on which those equations
hold to get all equations on one neighbourhood. These holomorphic units
make their zero germs equal.
Equality of finite cardinalities follows from that actual bijection.
For a common point $q$, translate the source by $q$ before this argument;
its derivative is the identity and the original inverse derivative
remains surjective. Dimension zero yields two empty selected coordinate
sets and is included.
\end{proof}

\begin{corollary}[Original manifold branches and intersections agree on overlaps]
\label{hir:original-normal-crossing-branch-overlap}
\uses{hir:actual-normal-crossing-chart-branch-correspondence,hir:general-nc-actual-local-snc-branches}
\lean{Hironaka.HironakaNormalCrossingsBranchOverlap.exists_unique_original_branch_equiv,Hironaka.HironakaNormalCrossingsBranchOverlap.exists_original_branch_equiv_intersections}
\leanok
At a point of an actual general normal-crossing carrier, any two actual
adapted flattenings have uniquely matching original coordinate branch
germs. Their local branch counts agree. Under the produced matching,
every finite intersection of original branches agrees as a germ.
\end{corollary}
\begin{proof}
Write both flattenings after the actual chart at the point, and translate
that chart centre to zero. Translation and each actual local inverse
give centred local holomorphic equivalences. The actual chart and its
inverse identify both selected zero unions with the original carrier
germ. Apply the preceding theorem. Transport each matched zero-set germ
through the actual chart back to the original manifold; continuity and
the local inverse identities preserve neighbourhood-germ equality.
For a finite intersection, intersect the finitely many neighbourhoods
on which the matched branch membership tests agree. On this common
neighbourhood all tests, and hence their conjunctions, agree. Equality
of the actual finite branch counts is already supplied by the produced
bijection, independently of the chosen flattening.
\end{proof}

\begin{remark}
This intrinsic correspondence does not imply that an ordered resolution
algorithm ignores initial labels. An invariant using a colexicographic
tie-break on ordered boundary indices can choose different centres under
a branch permutation. Canonical descent requires an invariant or a
first-blowup construction whose dependence on original unordered branches
has itself been proved, followed by genuine compatibility of subsequent
blowups. The correspondence above supplies the original overlap geometry.
\end{remark}
